\chapter{A Compilation, End to End}\label{ch:walkthrough}
% Inserted after the Backends chapter; chapter numbering is label-based, so
% downstream chapters renumber automatically. Content cycle: C10.
% Part 1 (this file, sections up to the EventList) traces one program from
% source through every IR stage to the per-worker event lists. Part 2 (the
% backend-lowering sections) shows the same event lists rendered to three
% backends.

The previous chapters described the language, the pipeline, and the backends one
concern at a time. This chapter puts them together. It takes a single program ---
the $3\times3$ stencil of \cref{lst:algo} under the distributed schedule of
\cref{lst:sched}, already familiar from \cref{sec:lang-example} --- and follows it
through every stage of the compiler, from the two source files to the per-worker
event lists and then, in the second half, to generated code on three different
backends. Nothing here is new machinery; everything is a concrete instance of a
pass defined in \cref{ch:architecture}. The purpose is to make the abstract
pipeline legible by watching one program move through it, with the actual
intermediate forms shown at each step.

The stencil is a good specimen because it is small enough to draw in full yet
exercises the parts of the compiler that do real work: it has a neighbourhood
access pattern that forces halo inference, an outer loop that is partitioned across
workers, an inner loop that is both tiled and reuse-optimised, and two transfers
with different IO semantics. Every figure and every intermediate form in this
chapter was produced by the shipping compiler on this exact input; the generated
code in the second half is reproduced from its output verbatim.

\section{The two inputs}\label{sec:wt-inputs}

Recall the division of labour. The algorithm file (\cref{lst:algo}) declares two
$16\times16$ arrays of 32-bit integers, \texttt{img\_in} and \texttt{img\_out};
three kernels --- a side-effecting (\emph{effectful}) \texttt{load\_image}, a pure
\texttt{blur3} taking nine neighbours, and an effectful \texttt{save\_image} (the
purity distinction is made precise in \cref{sec:wt-semantics}); and a dataflow body that
loads the image, writes each interior \texttt{img\_out[y][x]} from the nine
\texttt{img\_in} neighbours around it, and saves the result. It says nothing about
workers, transfers, tiling, or buffering.

The distributed schedule (\cref{lst:sched}) supplies exactly those things, and only
those things. It declares five workers --- a \texttt{host} and four compute workers
\texttt{w0}--\texttt{w3}; places \texttt{load\_image} and \texttt{save\_image} on
the host and \texttt{blur3} on the four compute workers; partitions the outer
\texttt{y} loop into row bands across them; tiles and reuse-optimises the inner
\texttt{x} loop; and asks that \texttt{img\_in} be streamed asynchronously with a
two-deep, event-notified buffer while \texttt{img\_out} is written back
synchronously. The compiler's task is to reconcile these two files into one
deterministic, sound, per-worker execution.

\section{The algorithm sublanguage, precisely}\label{sec:wt-semantics}

Before tracing the passes, it is worth stating precisely what the algorithm means,
because several later passes are sound only because of these semantics
(\cref{sec:lang-algorithm} introduced them; here they are pinned down on the
example).

\paragraph{Shaped, statically-sized arrays.} Every array has a shape fixed at
compile time. The constants \texttt{H} and \texttt{W} are evaluated to $16$ during
lowering, so \texttt{img\_in} and \texttt{img\_out} have fully concrete type
\texttt{i32[16][16]}. There are no runtime-sized arrays, and every index expression
can therefore be reasoned about as an affine function of the loop variables.

\paragraph{Shape-typed kernels with declared purity.} A kernel's signature gives
the shapes of its arguments and result, and a purity. \texttt{blur3} is
\emph{pure}: its result depends only on its arguments, so the compiler is free to
reorder, duplicate, or eliminate its calls. The reuse pass relies on the
elimination licence --- carrying a previously-read value over rather than re-reading
it --- and the partition pass relies on replication: because a pure kernel has no
hidden shared state, each worker may run it independently over its own band.
\texttt{load\_image} and \texttt{save\_image} are
\emph{effectful}: their relative order within a basic block is preserved and they
are never duplicated, because they touch the outside world. The kernel bodies live
in a sibling Rust file and are checked against these signatures before scheduling
begins; the algorithm sublanguage itself sees only the signatures.

\paragraph{Single assignment.} Each interior \texttt{img\_out[y][x]} is written
exactly once, and the lowering rejects a program that assigns the same datum twice
in one scope. Single assignment makes each datum's producer \emph{unique}; together
with the affine, static indices described below, that producer is also statically
\emph{identifiable}. At link time the compiler uses the coarser of these two facts
--- which kernel is placed on which worker --- to compute, for every array, its
producing and consuming worker sets; the partition pass later refines this to which
worker owns which cell.

\paragraph{Affine loops and the static firing model.} The loop bounds
(\texttt{1..H-1}) and every array index (\texttt{y-1}, \texttt{x+1}, and so on) are
affine in the loop variables with compile-time-constant coefficients. There is no
data-dependent control flow and no data-dependent trip count. As a consequence the
entire \emph{firing schedule} --- which kernel instances run, in what order, reading
and writing which array cells --- is determined statically, before any value is
known. This is the property that the entire middle-end leans on: because the firing
order is fixed by the program text rather than by runtime data, the compiler can
reason about decomposition, transfers, and soundness once, at compile time. That
answer --- which kernel instances fire, what data moves where, and whether the
coordination is bounded and deadlock-free --- is a property of the program text and
therefore holds for every input, independent of the values being computed; it is a
claim about the schedule and its coordination, not about the numerical results. The
price is the restricted class
(\cref{sec:lang-limits}); the return is everything that follows in this chapter.

\section{Parsing and lowering to the algorithm IR}\label{sec:wt-algoir}

The two files are parsed independently into abstract syntax trees and then lowered
into typed intermediate representations. Lowering the algorithm performs the work
the semantics above describe: it evaluates the constants (\texttt{H}, \texttt{W}
$\rightarrow 16$), resolves every array and kernel to a fully concrete shape-typed
form, checks the single-assignment discipline, and checks that each effectful
statement calls an effectful kernel. The result is the algorithm IR: a map of
resolved constants, a map of resolved data symbols with their concrete types, a map
of resolved kernels with their signatures and purity, and the dataflow body as a
list of statements --- a load, a doubly-nested loop whose innermost statement writes
\texttt{img\_out} from nine \texttt{img\_in} reads, and a save. Crucially, the
algorithm IR still contains no notion of a worker, a transfer, or a buffer; it is
the \emph{what}, fully typed, and nothing else.

The schedule is lowered in parallel into the schedule IR, which resolves the worker
declarations, the placements, the loop directives (the \texttt{partition} on
\texttt{y}; the \texttt{block} and \texttt{reuse} on \texttt{x}), and the two
transfer policies. At this point the two IRs are independent objects that have
never referred to each other.

\section{Linking: binding the algorithm to the schedule}\label{sec:wt-link}

Linking is where the two halves first meet (\cref{sec:arch-link}). It produces the
linked IR, which holds both source IRs together with the cross-references the
compiler derives by reconciling them: the placement of each kernel onto its worker
set, and, for each array, its producing worker set and its consuming worker sets.
For the stencil this resolution is the first concrete sign of the decomposition.
\texttt{load\_image} and \texttt{save\_image} bind to \texttt{host}; \texttt{blur3}
binds to the set $\{\texttt{w0},\texttt{w1},\texttt{w2},\texttt{w3}\}$.
\texttt{img\_in} is therefore produced on the host and consumed on all four compute
workers; \texttt{img\_out} is produced on the compute workers and consumed on the
host. Those producer/consumer facts are precisely the data movements the later
transfer pass must synthesise, but linking only records them; it inserts nothing
yet. Linking also enforces completeness: every kernel must be placed, and every
cross-worker producer/consumer pair must be reconcilable, or linking fails with a
diagnostic.

\section{The ACFG and its transforms}\label{sec:wt-acfg}

From the linked IR the compiler builds the application control-flow graph
(\cref{sec:arch-acfg}), a tree whose nodes are operations, sequences, and loops
(\cref{fig:wt-acfg}). For the stencil the initial tree is a sequence of three
items: an operation firing \texttt{load\_image} on the host; a repeat over
\texttt{y} containing a repeat over \texttt{x} containing an operation that fires
\texttt{blur3}; and an operation firing \texttt{save\_image} on the host. Each
operation carries a small dataflow record naming the arrays it reads and writes and
the index expressions it uses --- the nine \texttt{img\_in} accesses and the one
\texttt{img\_out} write, recorded as affine index expressions, which the inference
passes will read.

% fig:wt-acfg --- the stencil's application control-flow graph (ACFG).
% Faithful to chapters/07b-compilation-walkthrough.tex sec:wt-acfg:
%   initial tree = Seq[ Op(load_image@host),
%                       Repeat y:1..15 [ Repeat x:1..15 [ Op(blur3@{w0..w3}) ] ],
%                       Op(save_image@host) ]
% with each Op carrying a dataflow record (arrays read/written, affine indices),
% and the schedule transforms tagged on the nodes they rewrite. Ranges half-open.
\begin{figure}[tbp]
  \centering
  % ===== Panel (a): the ACFG control tree =====
  \begin{tikzpicture}[
      font=\small,
      seq/.style={draw, rounded corners, fill=gray!15, minimum height=7mm,
                  inner sep=3pt, align=center},
      op/.style={draw, fill=blue!10, minimum height=7mm, inner sep=3pt, align=center},
      rep/.style={draw, fill=orange!14, minimum height=7mm, inner sep=3pt, align=center},
      tree/.style={-{Latex[length=1.6mm]}},
      tag/.style={font=\scriptsize\itshape, fill=yellow!18, draw=yellow!55!black,
                  rounded corners=1pt, inner sep=2pt, align=center},
      lead/.style={densely dotted, yellow!55!black, -{Latex[length=1.2mm]}},
    ]
    \node[seq] (seq) at (0,0) {\texttt{Seq}};
    \node[op]  (load) at (-3.4,-1.5) {\texttt{Op}\,/\,\texttt{load\_image}\\@\,\texttt{host}};
    \node[rep] (ry)   at (0,-1.5)    {\texttt{Repeat} $y:1..15$};
    \node[op]  (save) at (3.4,-1.5)  {\texttt{Op}\,/\,\texttt{save\_image}\\@\,\texttt{host}};
    \node[rep] (rx)   at (0,-3.0)    {\texttt{Repeat} $x:1..15$};
    \node[op]  (blur) at (0,-4.5)    {\texttt{Op}\,/\,\texttt{blur3}\\@\,$\{$\texttt{w0..w3}$\}$};
    \draw[tree] (seq) -- (load);
    \draw[tree] (seq) -- (ry);
    \draw[tree] (seq) -- (save);
    \draw[tree] (ry) -- (rx);
    \draw[tree] (rx) -- (blur);

    % dataflow records (the affine read/write sets) to the right of each op
    \node[align=left, font=\scriptsize, right=2mm of blur, text width=42mm] (df) {%
      \textbf{reads} \texttt{img\_in[y$-$1..y$+$1][x$-$1..x$+$1]}\\
      \quad (9 affine accesses)\\
      \textbf{writes} \texttt{img\_out[y][x]}};
    \draw[densely dotted] (blur) -- (df);
    \node[font=\scriptsize, below=1mm of load] (dfl) {\textbf{writes} \texttt{img\_in}};
    \node[font=\scriptsize, below=1mm of save] (dfs) {\textbf{reads} \texttt{img\_out}};

    % schedule-transform tags, stacked in the empty lower-left, with leaders
    \node[tag, text width=24mm] (tpar) at (-3.5,-3.1)
      {\texttt{partition=rows}\\$\rightarrow$ 4 row bands};
    \node[tag, text width=24mm] (tblk) at (-3.5,-4.6)
      {\texttt{block=64} $+$ \texttt{reuse}\\$\rightarrow$ tiled, 3-slot window};
    \node[tag, text width=22mm] (thalo) at (6.7,-4.9)
      {halo pass\\$\rightarrow$ width-1 halo on \texttt{img\_in}};
    \draw[lead] (tpar.east) to[out=0,in=-135] (ry.south west);
    \draw[lead] (tblk.east) to[out=0,in=180] (rx.west);
    \draw[lead] (thalo.north) to[out=90,in=0] (df.south east);

    \node[font=\footnotesize\bfseries] at (0,0.95)
      {(a)\ \ ACFG control tree (transforms tagged in yellow)};
  \end{tikzpicture}

  \vspace{3mm}

  {\footnotesize\bfseries (b)\ \ dataflow DAG (the \emph{what}, before any decomposition)}\par
  \vspace{1.5mm}

  % ===== Panel (b): the dataflow DAG (resized to fit the text width) =====
  \resizebox{\textwidth}{!}{%
  \begin{tikzpicture}[
      font=\small,
      data/.style={draw, ellipse, fill=green!14, inner sep=2pt, align=center},
      kern/.style={draw, rounded corners=2pt, inner sep=3pt, align=center},
      flow/.style={-{Latex[length=2mm]}, thick},
    ]
    \node[kern, fill=blue!10] (k1) at (0,0) {\texttt{load\_image}\\\scriptsize effectful};
    \node[data] (d1) at (3.2,0) {\texttt{img\_in}\\\scriptsize\texttt{i32[16][16]}};
    \node[kern, fill=green!12] (k2) at (6.4,0) {\texttt{blur3}\\\scriptsize pure};
    \node[data] (d2) at (9.6,0) {\texttt{img\_out}\\\scriptsize\texttt{i32[16][16]}};
    \node[kern, fill=blue!10] (k3) at (12.6,0) {\texttt{save\_image}\\\scriptsize effectful};
    \draw[flow] (k1) -- (d1);
    \draw[flow] (d1) -- node[above, font=\scriptsize]{9 reads} (k2);
    \draw[flow] (k2) -- node[above, font=\scriptsize]{1 write/cell} (d2);
    \draw[flow] (d2) -- (k3);
  \end{tikzpicture}}

  \caption[The stencil's ACFG and dataflow DAG]{The application control-flow
    graph the compiler builds for the stencil. \textbf{(a)} The control tree: a
    \texttt{Seq} of an effectful load, a doubly-nested \texttt{Repeat} firing the
    pure \texttt{blur3}, and an effectful save (\texttt{@} denotes placement, so
    \texttt{@\{w0..w3\}} means the operation runs on all four compute workers).
    Each operation carries a dataflow record naming the arrays it reads and
    writes as affine index expressions (dotted). The schedule's transforms (yellow) rewrite or annotate specific
    nodes: \texttt{partition=rows} cuts the $y$ repeat into four bands,
    \texttt{block=64} with \texttt{reuse} tiles the $x$ repeat and adds a
    three-slot delay line, and halo inference records a width-one halo on
    \texttt{img\_in}. \textbf{(b)} The same body as a dataflow DAG --- still the
    \emph{what}, with no worker, transfer, or buffer in sight. Loop ranges are
    half-open ($1..15$ denotes $1,\dots,14$).}
  \label{fig:wt-acfg}
\end{figure}


The schedule's transforms are then applied to this tree, in the fixed order the
pipeline defines, each pass either rewriting the tree or annotating it. Throughout,
loop and band ranges are written half-open: $a..b$ denotes the integers
$a, a{+}1, \dots, b{-}1$.

\begin{itemize}
  \item \textbf{Tiling.} The \texttt{block=64} directive on \texttt{x} strip-mines
        the inner loop into a tile loop over an outer tile index and an intra-tile
        loop. The interior \texttt{x} range is $1..15$, of length 14; since the
        block size 64 exceeds it, there is a single tile here. The structure is
        nonetheless materialised, so the same code path handles the partial-tile
        case that arises when the block does not divide the range.
  \item \textbf{Partitioning.} The \texttt{partition=rows} directive on \texttt{y}
        is consumed by the row-band pass, which overrides each compute worker's
        \texttt{y} range with its band. The outer range $1..15$ likewise has length
        14, which does not divide evenly by four; the compiler distributes the
        remainder with a floor-with-spillover policy, handing
        $\lceil 14/4 \rceil = 4$ rows to the first two workers and
        $\lfloor 14/4 \rfloor = 3$ to the last two, giving the contiguous bands
        $1..5$, $5..9$, $9..12$, $12..15$ (\cref{fig:wt-bands}). The partition passes for
        per-worker ranges, row bands, and two-dimensional grids are mutually
        exclusive alternatives selected by the directive; here only the row-band
        pass fires.
  \item \textbf{Halo inference.} The halo pass reads \texttt{blur3}'s access
        pattern and finds that each output row $y$ reads input rows $y-1$, $y$, and
        $y+1$ (and analogously $x-1, x, x+1$): offsets of $-1$, $0$, $+1$ in both
        dimensions, recognised because each index is affine in the loop variable.
        It records a halo of width one on \texttt{img\_in} along the partitioned
        \texttt{y} axis. This is the fact that makes a row-band partition correct:
        a worker computing rows $1..5$ needs input rows $0..6$ --- its band plus one
        halo row above and below.
  \item \textbf{Reuse inference.} The \texttt{reuse} directive on \texttt{x} is
        discharged by the reuse pass, which recognises that as \texttt{x} advances
        by one the three-wide window of \texttt{img\_in} reads shifts by one, so two
        of every three reads can be carried over. It annotates the loop with a
        three-slot circular delay line per input row. This is the optimisation
        earlier iterations of the system could not express; here it falls out of a
        one-word directive and the access pattern.
\end{itemize}

% fig:wt-bands --- the row-band partition with the floor-with-spillover remainder
% and the width-1 halo. Faithful to sec:wt-acfg / sec:wt-inject of
% chapters/07b-compilation-walkthrough.tex:
%   interior y range 1..15 (length 14), partitioned 4/4/3/3 across w0..w3:
%     w0 = 1..5,  w1 = 5..9,  w2 = 9..12,  w3 = 12..15   (output bands, half-open)
%   halo width 1 on img_in along y => each worker receives band +- 1 row:
%     w0 input 0..6,  w1 input 4..10,  w2 input 8..13,  w3 input 11..16
%   rows 0 and 15 are the never-written border (interior is 1..15).
\begin{figure}[tbp]
  \centering
  \begin{tikzpicture}[
      font=\small,
      cell/.style={draw, minimum width=20mm, minimum height=4.2mm,
                   inner sep=0pt, anchor=north},
      brace/.style={decorate, decoration={brace, amplitude=4pt}, thick},
    ]
    % geometry: row r drawn at y = -r*0.42 cm (row 0 at top).
    \def\rh{0.42}
    \foreach \r in {0,...,15} {
      \pgfmathsetmacro\yy{-\r*\rh}
      % default border colour
      \def\fillcol{gray!18}
      \ifnum\r>0 \ifnum\r<5 \def\fillcol{blue!22}\fi\fi      % w0 band rows 1..4
      \ifnum\r>4 \ifnum\r<9 \def\fillcol{orange!28}\fi\fi    % w1 band rows 5..8
      \ifnum\r>8 \ifnum\r<12 \def\fillcol{green!28}\fi\fi    % w2 band rows 9..11
      \ifnum\r>11 \ifnum\r<15 \def\fillcol{red!22}\fi\fi     % w3 band rows 12..14
      \node[cell, fill=\fillcol] (row\r) at (0,\yy) {\scriptsize row \r};
    }
    % left header
    \node[font=\footnotesize\bfseries, align=center] at (0,0.55)
      {\texttt{img\_out} (written)};
    \node[font=\scriptsize, left=8mm of row0, anchor=east] {border (not written)};
    \node[font=\scriptsize, left=8mm of row15, anchor=east] {border (not written)};

    % band labels (output) on the left
    \node[font=\scriptsize, blue!50!black, left=1mm of row2.west, anchor=east, align=right]
      {\textbf{w0}\\out $1..5$\\(4 rows)};
    \node[font=\scriptsize, orange!55!black, left=1mm of row6.west, anchor=east, align=right]
      {\textbf{w1}\\out $5..9$\\(4 rows)};
    \node[font=\scriptsize, green!45!black, left=1mm of row10.west, anchor=east, align=right]
      {\textbf{w2}\\out $9..12$\\(3 rows)};
    \node[font=\scriptsize, red!55!black, left=1mm of row13.west, anchor=east, align=right]
      {\textbf{w3}\\out $12..15$\\(3 rows)};

    % halo-input ranges as braces on the right.  Each worker receives its band
    % +-1 halo row.  Braces drawn at increasing x to show the overlaps (halo).
    % w0 input rows 0..5 (range 0..6)
    \draw[brace, blue!50!black] ([xshift=4mm]row0.north east) -- ([xshift=4mm]row5.south east)
      node[midway, right=5pt, font=\scriptsize, align=left, blue!50!black]
      {w0 in: rows $0..6$};
    % w1 input rows 4..9 (range 4..10)
    \draw[brace, orange!55!black] ([xshift=16mm]row4.north east) -- ([xshift=16mm]row9.south east)
      node[midway, right=5pt, font=\scriptsize, align=left, orange!55!black]
      {w1 in: rows $4..10$};
    % w2 input rows 8..12 (range 8..13)
    \draw[brace, green!45!black] ([xshift=28mm]row8.north east) -- ([xshift=28mm]row12.south east)
      node[midway, right=5pt, font=\scriptsize, align=left, green!45!black]
      {w2 in: rows $8..13$};
    % w3 input rows 11..15 (range 11..16)
    \draw[brace, red!55!black] ([xshift=40mm]row11.north east) -- ([xshift=40mm]row15.south east)
      node[midway, right=5pt, font=\scriptsize, align=left, red!55!black]
      {w3 in: rows $11..16$};

    \node[font=\footnotesize\bfseries, align=center] at (4.0,0.55)
      {\texttt{img\_in} received (band $+$ halo)};
  \end{tikzpicture}
  \caption[Row-band partition and halo of the stencil]{The row-band partition of
    the stencil's outer loop. The interior range $1..15$ (rows $1$ through $14$;
    rows $0$ and $15$ are the never-written border) has length $14$, which does
    not divide evenly by four. The floor-with-spillover policy hands
    $\lceil 14/4\rceil = 4$ rows to \texttt{w0} and \texttt{w1} and
    $\lfloor 14/4\rfloor = 3$ to \texttt{w2} and \texttt{w3}, giving the
    contiguous output bands $1..5$, $5..9$, $9..12$, $12..15$ (left). Halo
    inference adds one row on each side along $y$, so each worker receives a
    \emph{larger} slice of \texttt{img\_in} (right braces): \texttt{w0} computing
    $1..5$ receives input rows $0..6$, \texttt{w1} computing $5..9$ receives
    $4..10$, and so on. The brace overlaps are the shared halo rows. All ranges
    are half-open.}
  \label{fig:wt-bands}
\end{figure}


After these passes the ACFG is no longer a faithful echo of the source loop nest:
it encodes four per-worker bands, a tiled and reuse-windowed inner loop, and the
knowledge that \texttt{img\_in} carries a one-row halo --- yet still contains no
transfers or synchronisations. Those are injected next.

\section{Injecting synchronisation and transfers}\label{sec:wt-inject}

Two final middle-end passes turn the decomposition implied by the ACFG into
explicit coordination: the synchronisation pass and the transfer pass
(\cref{sec:arch-transfer}). The synchronisation pass inserts the barriers that
bracket the distributed phase: one before the compute workers begin, after the host
has the input ready, and one after they finish, before the host collects the
output. Each barrier names its participants --- here all five workers --- and is
given a stable \emph{sync tag}, a compile-time identifier shared by every
participant of that barrier (\cref{sec:arch-petri}), so the projection can match the
same barrier across workers without a global walk.

The transfer pass then consumes the producer/consumer facts from linking and the
halo width from inference, and synthesises matched send/receive pairs. For
\texttt{img\_in} it emits, for each compute worker, a transfer from the host
carrying that worker's band \emph{plus its halo rows}: worker \texttt{w0}, which
computes rows $1..5$, receives input rows $0..6$; worker \texttt{w1}, computing
$5..9$, receives rows $4..10$; and so on. These transfers carry the schedule's
\texttt{img\_in} policy --- asynchronous, two-deep buffer, event notification. For
\texttt{img\_out} it emits, for each compute worker, a transfer back to the host
carrying just that worker's band of results, under the synchronous policy (a
single-deep, blocking channel: the host receives each band before the worker could
send another). The result is the fully elaborated ACFG: decomposition, halos, synchronisation, and IO
semantics all explicit, and still not one line of backend-specific code.

\section{The Petri net and the soundness gate}\label{sec:wt-net}

The elaborated ACFG is now lowered to a Petri net and checked
(\cref{sec:arch-gate,fig:wt-net}). In a Petri net, \emph{places} hold tokens up to a
capacity and \emph{transitions} fire by consuming tokens from their input places and
producing tokens in their output places. Here places model the things that hold
state --- each array region a worker owns, each transfer channel, and each barrier
--- and transitions model the things that happen: the host's \texttt{load\_image} firing,
the four \texttt{img\_in} sends and their matching receives, each worker's
\texttt{blur3} firings, the two barriers, the four \texttt{img\_out} sends and
receives, and the host's \texttt{save\_image} firing. A channel place has a finite
capacity drawn from the transfer's buffer policy --- two for the asynchronous
\texttt{img\_in} channels, one for the synchronous \texttt{img\_out} channels.

The gate runs its three checks in a fixed order (\cref{sec:arch-gate}). It first
confirms \emph{conflict-freedom}: no place offers a token to two competing
consumers. This is the load-bearing precondition. For the restricted net class this
language produces --- which by construction has no data-dependent token routing ---
conflict-freedom removes the only remaining source of choice, so every firing order
is a permutation of one partial order and reaches the same states; that is what
makes a single replayed order representative of all of them. The argument is
specific to this class and is not a general Petri-net property. On that basis the
gate derives one deterministic firing order and replays it, checking
\emph{boundedness} --- no channel place ever exceeds the capacity its buffer policy
sets --- and \emph{deadlock-freedom} --- no required transition is left permanently
un-fireable, whether because an input place never receives its token or because a
bounded output place stays full in a cyclic wait. For the stencil the net passes:
the two-deep input buffers never overflow, because the host pushes exactly one band
into each per-worker channel before reaching the barrier, well within the two-deep
capacity; and the barriers admit a completing order. Had the schedule
asked for a single-deep buffer where two were needed, or arranged a cyclic wait,
the gate would have rejected the program here, at compile time, naming the
offending place or transition --- not produced code that hangs. This check runs on
this program, as on every program, before any code is emitted.

% fig:wt-net --- the Petri net the elaborated ACFG lowers to, and the soundness
% gate. Faithful to sec:wt-net of chapters/07b-compilation-walkthrough.tex:
%   places  = array regions a worker owns, transfer channels, barriers;
%   transitions = host load_image, 4 img_in sends + 4 matching receives,
%                 each worker's blur3 firings, the two barriers,
%                 4 img_out sends + receives, host save_image.
%   channel capacities: 2 for the async img_in channels, 1 for the sync img_out.
% The four compute workers are a 4-fold-symmetric column; drawn once inside a
% plate ("x4") for legibility, with the host spine across the top. The two
% barriers are 5-way (host + w0..w3).
\begin{figure}[tbp]
  \centering
  \resizebox{\textwidth}{!}{%
  \begin{tikzpicture}[
      font=\small,
      place/.style={draw, circle, minimum size=9mm, inner sep=1pt, align=center,
                    fill=green!10},
      chan/.style={draw, circle, minimum size=10mm, inner sep=1pt, align=center,
                   fill=blue!10, very thick},
      tr/.style={draw, fill=gray!75, minimum width=2.2mm, minimum height=10mm,
                 inner sep=0pt},
      bar/.style={draw, fill=orange!70, minimum width=2.6mm, minimum height=34mm,
                  inner sep=0pt},
      a/.style={-{Latex[length=1.8mm]}},
      lbl/.style={font=\scriptsize, align=center},
    ]
    % ---- host spine (y = 2.6) ----
    \def\hy{2.6}
    \node[tr] (tload) at (0,\hy) {};   \node[lbl, above=0.5mm of tload] {\texttt{load\_image}};
    \node[place] (pin) at (1.5,\hy) {\texttt{img\_in}\\\scriptsize @host};
    \node[tr] (tsend) at (3.0,\hy) {}; \node[lbl, above=0.5mm of tsend] {send};
    \node[tr] (trout) at (12.6,\hy) {};\node[lbl, below=0.5mm of trout] {recv};
    \node[place] (pout) at (14.1,\hy) {\texttt{img\_out}\\\scriptsize @host};
    \node[tr] (tsave) at (15.6,\hy) {};\node[lbl, above=0.5mm of tsave] {\texttt{save\_image}};

    % ---- a single worker column (y = 0), enclosed by a plate ----
    \node[chan] (cin) at (4.5,1.3) {cap\,2};
    \node[lbl, left=0.5mm of cin] {\texttt{img\_in}\\chan};
    \node[tr] (trin) at (6.0,0) {};   \node[lbl, below=0.5mm of trin] {recv};
    \node[place] (pband) at (7.4,0) {band\\\scriptsize $+$halo};
    \node[tr] (tblur) at (8.8,0) {}; \node[lbl, below=0.5mm of tblur] {\texttt{blur3}};
    \node[place] (pob) at (10.0,0) {out\\\scriptsize band};
    \node[tr] (tso) at (11.1,0) {};  \node[lbl, below=0.5mm of tso] {send};
    \node[chan] (cout) at (12.6,1.3) {cap\,1};
    \node[lbl, right=0.5mm of cout] {\texttt{img\_out}\\chan};

    % plate around the worker column
    \begin{scope}[on background layer]
      \node[draw, dashed, rounded corners, fill=yellow!8, inner sep=4mm,
            fit=(cin)(trin)(pband)(tblur)(pob)(tso)(cout),
            label={[font=\footnotesize\bfseries, fill=yellow!18, rounded corners,
                    inner sep=2pt, yshift=-2mm]below:$\times\,4$ --- one column per
                    worker (\texttt{w0}--\texttt{w3}): bands $1..5$ / $5..9$ /
                    $9..12$ / $12..15$}]
            (plate) {};
    \end{scope}

    % ---- the two 5-way barriers ----
    \node[bar] (b1) at (5.0,1.4) {};
    \node[lbl, orange!60!black] at (5.0,3.7) {barrier 1\\\scriptsize (pre-compute, 5-way)};
    \node[bar] (b2) at (11.9,1.4) {};
    \node[lbl, orange!60!black] at (11.9,3.7) {barrier 2\\\scriptsize (post-compute, 5-way)};

    % ---- arcs ----
    \draw[a] (tload) -- (pin);
    \draw[a] (pin) -- (tsend);
    \draw[a] (tsend) -- (cin);            % host pushes into the img_in channel
    \draw[a] (cin) -- (trin);             % worker receives
    \draw[a] (trin) -- (pband);
    \draw[a] (pband) -- (tblur);
    \draw[a] (tblur) -- (pob);
    \draw[a] (pob) -- (tso);
    \draw[a] (tso) -- (cout);             % worker pushes into img_out channel
    \draw[a] (cout) -- (trout);           % host receives
    \draw[a] (trout) -- (pout);
    \draw[a] (pout) -- (tsave);
    % barrier synchronisation arcs (host + worker both meet each barrier)
    \draw[a, orange!60!black] (tsend) -- (b1);
    \draw[a, orange!60!black] (b1) -- (trin);
    \draw[a, orange!60!black] (tso) -- (b2);
    \draw[a, orange!60!black] (b2) -- (trout);
    % host participates in BOTH barriers: it passes barrier 1, idles, then meets
    % barrier 2 (the back-to-back "barrier; barrier" in the host's event list).
    \draw[a, orange!60!black] (b1.north) to[out=70,in=110] (b2.north);
  \end{tikzpicture}}
  \caption[The stencil's Petri net and soundness gate]{The Petri net the
    elaborated ACFG lowers to. \emph{Places} (circles) model state --- the array
    regions a worker owns and the transfer channels --- and \emph{transitions}
    (bars) model events: the host's \texttt{load\_image}, the four
    \texttt{img\_in} sends and matching receives, each worker's \texttt{blur3}
    firings, the four \texttt{img\_out} sends and receives, and the host's
    \texttt{save\_image}. The four compute workers form a 4-fold-symmetric
    column, drawn once inside the dashed plate. Channel places (marked
    ``cap'') carry the capacity their buffer policy sets: \textbf{2} for the
    asynchronous \texttt{img\_in} channels, \textbf{1} for the synchronous
    \texttt{img\_out} channels. The two orange bars are the 5-way (host $+$
    \texttt{w0}--\texttt{w3}) pre- and post-compute barriers; the orange arcs
    are their barrier-place arcs (no participant proceeds until all five have
    arrived), and the arc arching over the top is the host idling from the
    first barrier to the second. The gate checks, in order, conflict-freedom
    (no place offers a token to two competing consumers), then --- on one
    replayed firing order --- boundedness (no \texttt{img\_in} channel ever
    exceeds 2) and deadlock-freedom, before any code is emitted.}
  \label{fig:wt-net}
\end{figure}


\section{Projection to per-worker event lists}\label{sec:wt-events}

The last middle-end step projects the elaborated ACFG into one totally-ordered
event list per worker (\cref{sec:arch-projection,fig:wt-events}). The walk is
deterministic, so the lists are reproducible byte-for-byte across compilations. The
events are drawn from a small fixed vocabulary --- fire a kernel, push a data region
to a worker, wait for a region from a worker, participate in a barrier, and loop ---
and they are the \emph{sole} interface to the backend: everything the decomposition
decided is encoded in them, and nothing else crosses the line.

For the stencil the host's list reads: fire \texttt{load\_image}; push
\texttt{img\_in} to \texttt{w0}, \texttt{w1}, \texttt{w2}, \texttt{w3}; barrier;
barrier; wait for \texttt{img\_out} from \texttt{w0}, \texttt{w1}, \texttt{w2},
\texttt{w3}; fire \texttt{save\_image}. The host stages all four input bands into
the buffered channels, then the two barriers bracket the phase during which it has
nothing to do, then it collects the four output bands. The two barriers are the
same pre- and post-compute synchronisations of \cref{sec:wt-inject}: the first
releases the workers once every input band has been pushed, the second waits until
every worker has finished and pushed its output. Each compute worker's list is the
mirror image. Worker \texttt{w0}'s reads: barrier; wait for \texttt{img\_in} from
the host; a loop over its band rows $1..5$, whose body is the tiled, reuse-windowed
inner loop firing \texttt{blur3}; push \texttt{img\_out} to the host; barrier. The
loop event carries the iteration structure rather than an unrolled body, so a
backend emits a loop construct rather than one replicated body per iteration,
keeping generated code compact regardless of trip count.

% fig:wt-events --- the per-worker event lists as a swimlane timeline.
% Faithful to sec:wt-events of chapters/07b-compilation-walkthrough.tex:
%   host : fire load_image; push img_in to w0,w1,w2,w3; barrier; barrier;
%          wait img_out from w0,w1,w2,w3; fire save_image.
%   w_i  : barrier; wait img_in from host; loop over band { tiled+reuse blur3 };
%          push img_out to host; barrier.
% Time flows downward. The two barriers are horizontal 5-way sync lines; the host
% idles (empty lane) between them while the workers compute -- which is exactly
% the host's back-to-back "barrier; barrier" in its list.
\begin{figure}[tbp]
  \centering
  \resizebox{\textwidth}{!}{%
  \begin{tikzpicture}[
      font=\small,
      fire/.style={draw, rounded corners=3pt, fill=blue!12, minimum height=6mm,
                   inner sep=3pt, align=center},
      io/.style={draw, fill=green!12, minimum height=6mm, inner sep=3pt, align=center},
      loop/.style={draw, fill=orange!14, minimum height=13mm, inner sep=3pt,
                   align=center, font=\scriptsize},
      bar/.style={very thick, orange!65!black},
      lane/.style={font=\footnotesize\bfseries},
      push/.style={-{Latex[length=1.8mm]}, blue!45!black, thick},
      ret/.style={-{Latex[length=1.8mm]}, green!40!black, thick},
    ]
    % lane x positions
    \def\xh{0} \def\xa{3.4} \def\xb{5.5} \def\xc{7.6} \def\xd{9.7}
    % lane headers
    \node[lane] at (\xh,0.7) {\texttt{host}};
    \node[lane] at (\xa,0.7) {\texttt{w0}};
    \node[lane] at (\xb,0.7) {\texttt{w1}};
    \node[lane] at (\xc,0.7) {\texttt{w2}};
    \node[lane] at (\xd,0.7) {\texttt{w3}};

    % --- host lane above barrier 1 ---
    \node[fire] (hload) at (\xh,0)   {fire \texttt{load\_image}};
    \node[io]   (hpush) at (\xh,-1.1){push \texttt{img\_in}\\$\rightarrow$ w0,w1,w2,w3};

    % --- barrier 1 (5-way) ---
    \def\byone{-2.0}
    \draw[bar] (-1.2,\byone) -- (10.9,\byone);
    \node[font=\scriptsize, anchor=west, orange!65!black] at (10.95,\byone) {barrier 1};

    % --- worker lanes between the barriers ---
    \foreach \x/\name/\band in {\xa/w0/{1..5}, \xb/w1/{5..9}, \xc/w2/{9..12}, \xd/w3/{12..15}} {
      \node[io]   (\name wait) at (\x,-2.75) {wait\\\texttt{img\_in}};
      \node[loop] (\name loop) at (\x,-4.15) {loop $y:\band$\\\scriptsize tiled $+$ reuse\\\scriptsize fire \texttt{blur3}};
      \node[io]   (\name out)  at (\x,-5.55) {push\\\texttt{img\_out}};
    }
    % host idles between the two barriers (empty lane) -- its list is "barrier; barrier"
    \node[font=\scriptsize\itshape, gray, align=center] at (\xh,-4.15)
      {(host idle:\\\texttt{barrier;}\\\texttt{barrier})};

    % --- barrier 2 (5-way) ---
    \def\bytwo{-6.4}
    \draw[bar] (-1.2,\bytwo) -- (10.9,\bytwo);
    \node[font=\scriptsize, anchor=west, orange!65!black] at (10.95,\bytwo) {barrier 2};

    % --- host lane below barrier 2 ---
    \node[io]   (hwait) at (\xh,-7.2) {wait \texttt{img\_out}\\$\leftarrow$ w0,w1,w2,w3};
    \node[fire] (hsave) at (\xh,-8.3) {fire \texttt{save\_image}};

    % --- intra-lane order edges (host) ---
    \draw[-{Latex[length=1.4mm]}] (hload) -- (hpush);
    \draw[-{Latex[length=1.4mm]}] (hwait) -- (hsave);
    % --- intra-lane order edges (workers) ---
    \foreach \name in {w0,w1,w2,w3} {
      \draw[-{Latex[length=1.4mm]}] (\name wait) -- (\name loop);
      \draw[-{Latex[length=1.4mm]}] (\name loop) -- (\name out);
    }
    % --- cross-lane transfers: img_in pushes (host -> workers) ---
    \foreach \name in {w0,w1,w2,w3} {
      \draw[push] (hpush.east) to[out=0,in=120] (\name wait.north);
    }
    % --- cross-lane transfers: img_out returns (workers -> host) ---
    \foreach \name in {w0,w1,w2,w3} {
      \draw[ret] (\name out.south) to[out=-120,in=0] (hwait.east);
    }

    % --- legend ---
    \begin{scope}[shift={(-1.0,-9.4)}]
      \node[fire, minimum height=4mm] (lf) at (0,0) {fire};
      \node[io, minimum height=4mm, right=2mm of lf] (li) {push/wait};
      \node[loop, minimum height=4mm, right=2mm of li] (ll) {loop};
      \draw[bar] ([xshift=4mm]ll.east) -- ++(0.9,0);
      \node[font=\scriptsize, right=1mm of ll, xshift=10mm] {barrier (5-way sync)};
    \end{scope}
  \end{tikzpicture}}
  \caption[Per-worker event lists of the stencil]{The per-worker event lists, the
    sole interface between the middle end and any backend, as a swimlane
    timeline (time flows down). The host fires \texttt{load\_image}, pushes one
    \texttt{img\_in} band into each worker's buffered channel, meets both
    barriers, collects the four \texttt{img\_out} bands, and fires
    \texttt{save\_image}. Each worker meets barrier 1, waits for its
    \texttt{img\_in} band, runs the tiled, reuse-windowed \texttt{blur3} loop over
    its row band, pushes its \texttt{img\_out} band back, and meets barrier 2. The
    two horizontal bars are the 5-way pre- and post-compute barriers; between them
    the host lane is empty --- exactly the back-to-back \texttt{barrier; barrier}
    in the host's list --- while the workers compute. Blue arrows are
    \texttt{img\_in} transfers (host to worker), green arrows are
    \texttt{img\_out} transfers (worker to host); the \texttt{img\_in} pushes
    cross barrier 1 because the channels are buffered (depth 2). The loop event
    carries the iteration structure, not an unrolled body.}
  \label{fig:wt-events}
\end{figure}


These event lists are the output of the middle end and the input to every backend.
The same host-and-four-workers lists just described are what the next sections lower
to shared-memory threads, to MPI ranks, and to bare-metal microcontrollers --- the
same events, three very different renderings.

\section{One event list, three renderings}\label{sec:wt-lowering}

A backend consumes the per-worker event lists and nothing else
(\cref{sec:arch-codegen}). It renders each event into code for one combination of
runtime and transport, and the central observation of this chapter is visible only
once that rendering is laid side by side across backends: the events that compute
--- fire a kernel, loop --- go through a \emph{shared} renderer, while only the events
that communicate --- push, wait, barrier --- are rendered per-backend
(\cref{fig:wt-lowering}). This is the middle-end/presentation boundary
(\cref{sec:arch-boundaries}) --- the line separating the shared, backend-independent
lowering from the per-backend rendering of transport --- made mechanical: there is one
emitter for the inside of a worker and a small per-backend emitter for the seams
between workers.

The consequence is literal, not approximate. Emitting the stencil's distributed event
lists to the shared-memory \texttt{pthreads-async} backend and to the
distributed-memory \texttt{mpi-nonblocking} backend produces, for each compute
worker, \emph{byte-for-byte the same compute body} --- the same tiled, reuse-windowed
loop, the same nine-argument \texttt{blur3} call, the same circular-buffer index
arithmetic, down to the character (the transport sections of the two files, of
course, differ). The two generated programs diverge only where a worker sends or
receives a band or meets a barrier. The sections below show each of the three
renderings of the same event lists from \cref{sec:wt-events} --- two of them, the
hosted backends, running that exact distributed schedule, the third running a
synchronous variant its transport requires --- and the last explains why all three
agree to the byte.

\section{Shared memory: the pthreads rendering}\label{sec:wt-pthreads}

The \texttt{pthreads-async} backend maps each worker to an operating-system thread
and each transfer channel to a bounded ring buffer shared between two threads
(\cref{sec:be-tier1}). The host's \texttt{main} allocates one ring per transfer ---
four for the \texttt{img\_in} sends, each a two-slot ring matching the schedule's
\texttt{buffer=2}, and four single-slot rings for the synchronous \texttt{img\_out}
returns --- and two five-party barriers, then spawns the four compute-worker threads
and runs the host's own event list inline.

Each event renders to a shared-memory primitive. A \emph{push} becomes a
\texttt{ring.push(payload)} into the destination worker's channel; a \emph{wait}
becomes a \texttt{ring.wait()} followed by a copy of the received band into the
worker's local array; a \emph{barrier} becomes a \texttt{barrier.wait()} on the
shared five-party barrier; and a \emph{fire} becomes a direct call to the kernel.
The ring is a condition-variable queue --- a producer blocks when it is full, a
consumer blocks when it is empty --- so the asynchronous, event-notified semantics
the schedule asked for are realised by the operating system's condition variables
rather than by any polling. The host stages the four input bands into the
two-deep rings, both barriers pass, and the four output bands are collected; the
compute worker for the first band runs exactly the loop \cref{sec:wt-acfg}
described, reading its received rows $0..6$ and writing its band rows $1..5$.

\section{Distributed memory: the MPI rendering}\label{sec:wt-mpi}

The \texttt{mpi-nonblocking} backend maps each worker to an MPI rank and emits a
single program that every rank runs, branching on its rank to execute that worker's
event list (\cref{sec:be-tier2,fig:mpi-spmd}). Where the shared-memory backend
spawned four threads from a host \texttt{main}, the MPI backend emits one \texttt{match
world.rank()} --- a branch on the process's rank within the MPI communicator --- with
five arms: rank $0$ is the host, ranks $1$ through $4$ are the compute workers, and
the program is launched as five processes with
\texttt{mpiexec~-n~5}. The same event list that drove the host thread now drives
rank $0$; the same list that drove the first compute thread now drives rank $1$.

The transport primitives are entirely different from the ring buffers, yet they
realise the same event semantics. A transfer channel becomes a small wrapper over an
MPI peer and a message tag (a small integer the receiver uses to select which
incoming message to take). A \emph{push} is a buffered non-blocking send
(\texttt{MPI\_Ibsend} into a process-attached send buffer) whose completion is local:
it returns once the payload has been copied into that buffer, without waiting for the
peer to post its receive. This is what gives the asynchronous, buffered behaviour the
schedule requested, and what keeps a wait-before-push exchange from deadlocking --- if
two workers each waited to receive before sending, neither would proceed, whereas a
local-completion send lets each post its send first and then wait. A \emph{wait} is a
matched probe (\texttt{MPI\_Mprobe}, which learns the incoming element count) followed
by a non-blocking matched receive and its completion. A \emph{barrier} is
\texttt{MPI\_Barrier} over the whole communicator. Each channel pins its MPI message
tag to the transfer's compile-time identity --- every distinct transfer edge in the
event list gets a unique tag --- so two concurrent transfers between the same pair of
ranks carry different tags and cannot be matched by the wrong receive. The rank assignment, the message tags, and the receive slices are all
computed by the same projection that produced the event lists; the backend only
chooses how to spell a push and a wait.

And the inside of each rank is the shared compute body unchanged. Rank $1$ receives
rows $0..6$ into the same array offset (\texttt{img\_in[0..96]}: six rows of sixteen
in row-major order), runs the byte-identical tiled, reuse-windowed \texttt{blur3}
loop, and sends its output array back to the host, which recovers this worker's band
rows $1..5$ at the same offset the pthreads host used (\texttt{img\_out[16..80]}) ---
the identical arithmetic the pthreads thread ran, wrapped in \texttt{MPI\_Ibsend} and
\texttt{MPI\_Mprobe} instead of \texttt{ring.push} and \texttt{ring.wait}.

\section{The capability boundary}\label{sec:wt-capability}

The two renderings above were both possible because both backends advertise the
capabilities the stencil's distributed schedule demands: asynchronous transfers, a
buffer of depth two, and event notification. Not every backend does, and the event
contract is where the mismatch is caught. The schedule's \texttt{img\_in} transfer
asks for \texttt{async}, \texttt{buffer=2}, and \texttt{notify=event}; a backend
whose capability declaration offers only synchronous, single-buffer, blocking
transfers cannot honour it. When the same program is aimed at such a backend ---
\texttt{mpi-blocking}, or the embedded backend of the next section --- compilation
\emph{fails}, with a diagnostic naming each unmet demand: that the transfer requests
asynchrony the backend does not support, a buffer depth beyond the backend's maximum
of one, and an event notification mode the backend does not offer. This is not a
runtime surprise or a silent downgrade to blocking behaviour --- which would make the
same schedule mean different things on different backends --- but a compile-time
rejection (\cref{sec:be-matrix}). A schedule and a backend compose only when the
backend's capabilities are a superset of what the event list demands; here a capable
sibling such as \texttt{pthreads-async} carries the schedule, and the synchronous
backends carry the synchronous schedules of the same algorithm instead.

\section{Bare metal: the embedded rendering}\label{sec:wt-embedded}

The embedded backend renders an event list to a \texttt{no\_std} crate for a
microcontroller, with no operating system, no heap, and no threads
(\cref{sec:be-tier3}). Because its transport is a synchronous UART or DMA fabric, it
declares only synchronous, single-buffer, blocking transfers, so --- as just shown ---
it rejects the stencil's asynchronous distributed schedule. The same stencil
\emph{algorithm} under a synchronous schedule --- one whose transfers drop the
\texttt{async}, \texttt{buffer}, and event-notify directives in favour of plain
blocking transfers --- lowers cleanly, and exhibits the third rendering of the event
contract.

Three things change relative to the hosted backends, and all three follow from the
target rather than from the algorithm. First, storage: the heap-allocated
\texttt{Vec} arrays of the hosted backends become fixed-size stack arrays
(\texttt{[i32; 256]}), since there is no allocator. Second, the effectful kernels:
where the hosted backends call \texttt{load\_image} and \texttt{save\_image} as
ordinary functions, the embedded backend lowers them to calls on a
\emph{hardware-abstraction trait} --- a Rust interface that each per-microcontroller
shim implements concretely --- so loading the image becomes a region allocation and a
DMA wait, and saving it becomes a DMA push to a peripheral followed by a wait. Third,
the transport: a cross-worker \emph{push} becomes a \texttt{link\_push} on the trait,
keyed by the transfer's sequence tag (the same compile-time transfer identity the MPI
backend mapped to a message tag), and a \emph{wait} becomes a matching
\texttt{link\_recv}; a \emph{barrier} becomes an interrupt-driven synchronisation
hook. The generic backend emits against the trait once; each microcontroller family
supplies a shim that binds these hooks to real DMA descriptors, interrupt vectors,
and memory addresses.

The \emph{compute kernel}, though, is the same: the pure \texttt{blur3} body is
copied verbatim into the embedded crate, character for character. The loop that
calls it differs from the hosted renderings --- this is the synchronous naive
schedule, with no partitioning across workers and no reuse window, so the loop is an
unpartitioned doubly-nested sweep over the whole image rather than the tiled,
reuse-windowed per-band loop the distributed schedule produced. What does not change
is the arithmetic each \texttt{blur3} call performs: the only kernel that changes
\emph{form} is the effectful IO, because only it touches hardware. For a multi-worker embedded
schedule the cross-worker pushes and waits become real inter-microcontroller UART
transfers, and a boot order is computed so that each receiver enables its link
before its sender transmits, because a bare-metal UART has no buffer to hold an
early byte (\cref{sec:be-tier3}); the dedicated multi-microcontroller examples
exercise that path under co-simulation and reproduce their reference output
byte-exactly (\cref{sec:res-tiers}).

\section{Why the three agree to the byte}\label{sec:wt-byteident}

The three renderings just walked through differ in almost everything a systems
programmer would notice --- threads versus processes versus bare-metal cores; ring
buffers versus MPI messages versus UART links; a heap versus fixed arrays --- and yet
the hosted two produce byte-identical output, and the embedded one reproduces the
same reference bytes under co-simulation. The walkthrough makes the reason concrete
rather than asserted. For the two backends running the \emph{same} (distributed)
schedule, the compute that determines every output byte is emitted by one shared
single-worker renderer, so it is the \emph{same source text}: the \texttt{blur3}
calls and their reuse-buffer index arithmetic are character-for-character identical
between the shared-memory and MPI programs. The embedded backend runs a different
(synchronous) schedule and so emits a different loop, but the value-determining part
--- the \texttt{blur3} kernel body --- is the same verbatim kernel. In every case that
arithmetic is integer arithmetic, which is deterministic and reorder-independent by
the language's design (\cref{sec:lang-limits}), so no backend's scheduling choices
can perturb a value. What the backends \emph{do} differ in --- how a band is moved from
one worker to another --- is value-preserving: a ring, an \texttt{MPI\_Ibsend}, and a
\texttt{link\_push} each move the band as an opaque byte payload, none reinterpreting
or recomputing it, and the soundness gate has already established that none overflows
or deadlocks for this program. For this stencil, whose workers write disjoint output
bands with no cross-worker reduction, identical compute over a value-preserving
transport yields identical output --- where workers instead combine into a shared
accumulator, byte-identity additionally requires the combine to be order-independent,
which the integer-arithmetic restriction supplies (\cref{sec:res-bitidentity}). The
cross-backend differential test of \cref{ch:validation} is the
empirical check that this reasoning holds, and the walkthrough is why it holds.
